% Operating curves. Both panels read curve.dat / scale.dat, which
% scripts/paper_numbers.py writes out of results/, so a plotted point cannot
% drift from the run that produced it. Drafted to match the text: hairline
% axes, no boxes, no fills, marks a reader can count.
\begin{figure}[t]
\centering
\pgfplotsset{
  draft/.style={
    width=\columnwidth, height=41mm,
    axis lines=left, axis line style={line width=0.4pt},
    tick align=outside, tick style={line width=0.4pt, black},
    ymajorgrids, grid style={black!12, line width=0.3pt},
    label style={font=\scriptsize}, tick label style={font=\scriptsize},
    legend style={font=\scriptsize, draw=none, fill=none, cells={anchor=west},
                  inner sep=1pt},
    every axis plot/.append style={line width=0.7pt, mark size=1.7pt},
  }}
\begin{tikzpicture}
\begin{axis}[draft, xlabel={차단 반경 $r$}, ylabel={SC / 다양성 (LPIPS)},
             xmin=-0.03, xmax=1.03, ymin=0, ymax=0.8,
             legend pos=north east]
  \addplot[mark=square*, black] table[x=r, y=sc] {curve.dat};
  \addlegendentry{구조 순응도}
  \addplot[mark=o, black, densely dashed] table[x=r, y=div] {curve.dat};
  \addlegendentry{표본 다양성}
\end{axis}
\end{tikzpicture}

\vspace{2mm}

\begin{tikzpicture}
\begin{axis}[draft, height=32mm,
             xlabel={구조 순응도}, ylabel={FID},
             xmin=-0.03, xmax=0.8, ymin=45, ymax=112,
             legend columns=2,
             legend style={at={(0.5,-0.42)}, anchor=north}]
  % curve.dat의 행 순서는 r = 0, 0.05, 0.1, 0.2, 0.3, 0.5, 0.7, 1이며 1--4행이
  % 주장하는 구간이다. \thisrow{r} 표현식 필터가 pgfplots 파서를 통과하지 못해
  % 행 인덱스로 나눈다.
  \addplot[mark=square*, black, skip coords between index={0}{1}, skip coords between index={5}{8}] table[x=sc, y=fid] {curve.dat};
  \addlegendentry{차단 반경, 자기 계측기}
  \addplot[only marks, mark=square, black!45, skip coords between index={1}{5}] table[x=sc, y=fid] {curve.dat};
  \addlegendentry{차단 반경 $r > 0.3$}
  \addplot[mark=diamond*, black, densely dotted] coordinates {(\fixInstZeroFive,\fidMainZeroFive) (\fixInstOneZero,\fidMainOneZero) (\fixInstTwoZero,\fidMainTwoZero) (\fixInstThreeZero,\fidMainThreeZero)};
  \addlegendentry{차단 반경, 계측기 $r_m{=}0.1$}
  \addplot[mark=o, black, densely dashed] table[x=sc, y=fid] {scale.dat};
  \addlegendentry{무작위-$r$, $r{=}0.1$에서 스케일}
  \addplot[mark=triangle, black!55, dotted] coordinates {(\sscSCQuarter,\sscFidQuarter) (\sscSCHalf,\sscFidHalf) (\sscSCThreeQ,\sscFidThreeQ) (\sscSCOne,\sscFidOne)};
  \addlegendentry{고정-$r{=}0.1$ 어댑터, 스케일}
\end{axis}
\end{tikzpicture}
\caption{\textbf{위:} 차단 반경에 대한 구조 순응도와 표본 다양성(표~\ref{tab:main}의
값이며 SEM과 시드 sd도 거기에 있다). \textbf{아래:} 구조 순응도에 대한 FID, 점마다 FID 추정은
하나다. 채운 사각형은 $r \in [0.05, 0.3]$에서 자기 차단 반경으로 채점한 스윕, 빈 회색 사각형은
같은 스윕의 구간 밖이다. 마름모는 같은 생성물을 계측기 $r_m{=}0.1$로 다시 채점한 것이다. 원과
삼각형은 무작위-$r$ 어댑터와 고정-$r{=}0.1$ 어댑터를 $r{=}0.1$에서 스케일한 것으로, 이미
$r_m{=}0.1$을 쓴다. 그 한 계측기에서 차단 반경 $r{=}0.3$은 FID \fidMainThreeZero{}에서
\fixInstThreeZero{}에, 세기 0.50은 FID \fidHalfScale{}에서 \fixInstHalfScale{}에 이른다.
두 패널 모두 평가 출력에서 직접 그렸다.}
\label{fig:curve}
\end{figure}
